
\subsubsection{Comparaison des méthodologies courantes}

Le tableau suivant (Tableau :~\ref{tab:comparaison_methodologies}) compare les méthodologies les plus couramment utilisées dans le développement logiciel, en mettant en évidence leurs caractéristiques clés :

\begin{table}[H]
    \centering
    \begin{tabular}{|m{3cm}|c|c|c|c|}
        \hline
        \textbf{Caractéristiques} & \textbf{Scrum} & \textbf{Waterfall} & \textbf{Lean} & \textbf{Kanban} \\
        \hline
        \textbf{Adaptabilité aux changements} & \checkmark & X & X & \checkmark \\
        \hline
        \textbf{Livraisons fréquentes} & \checkmark & X & X & X \\
        \hline
        \textbf{Collaboration renforcée} & \checkmark & X & X & \checkmark \\
        \hline
        \textbf{Réponse rapide aux changements} & \checkmark & X & X & X \\
        \hline
        \textbf{Gestion des risques} & \checkmark & \checkmark & \checkmark & \checkmark \\
        \hline
        \textbf{Niveau d'implication du client} & Élevé & Faible & Moyen & Moyen \\
        \hline
        \textbf{Efficacité des coûts} & Élevée & Moyenne & Élevée & Moyenne \\
        \hline
        \textbf{Assurance qualité} & Continue & Finale & Continue & Continue \\
        \hline
        \textbf{Scalabilité} & Moyenne & Élevée & Élevée & Moyenne \\
        \hline
        \textbf{Maintenance et support} & Facile & Difficile & Facile & Facile \\
        \hline
    \end{tabular}
    \caption{Comparaison des méthodologies de développement logiciel}
    \label{tab:comparaison_methodologies}
\end{table}
\subsubsection{Conclusion}

% Comme le montre le tableau, Scrum se distingue par sa flexibilité, son approche itérative et sa capacité à s'adapter aux exigences changeantes. Ces caractéristiques sont essentielles pour le développement d'un système complexe comme celui de la collecte de données et de la gestion des alertes en salle de réveil. 
Par conséquent, Scrum est la méthodologie la plus adaptée pour ce projet.
Les principales raisons du choix de Scrum sont les suivantes :
\begin{itemize}
    \item \textbf{Adaptabilité aux changements} : Les besoins des utilisateurs et les exigences techniques peuvent évoluer au fil du projet, et Scrum permet d’y répondre efficacement.
    \item \textbf{Développement itératif} : Chaque sprint permet d’obtenir un incrément fonctionnel du produit, facilitant ainsi les tests et l’amélioration continue.
    \item \textbf{Collaboration et transparence} : Les réunions régulières favorisent la communication entre les équipes et les parties prenantes, assurant une meilleure gestion du projet.
    \item \textbf{Amélioration continue} : L’analyse régulière des performances et des difficultés permet une optimisation constante du processus de développement.
\end{itemize}